\section{张量是多线性映射，也是多维数组}
\label{sec:03-Linear-Algebra/09-Tensors/01}

一段彩色视频在计算机里长这样：

\[
X\in\R^{T\times H\times W\times C},
\]

四个索引分别管时间、画面高度、画面宽度和颜色通道。你当然可以把后三个索引合并，拍平成一张 \(T\times (HWC)\) 的矩阵——每一行是一帧，每一列是一个像素-通道组合。也可以换一种拍法：把时间和高度合并，得到 \((TH)\times (WC)\) 的矩阵。数值一个没丢，但同一个数组拆成了完全不同的矩阵形状。

这制造了一个认知上的困惑。如果你把同一组数字摊平成一张表是矩阵 A，换个摊法又成了矩阵 B，那它们到底表示同一个数学对象还是两个不同的东西？计算接口把一切多于两个索引的数组都叫 tensor，形状对了就能乘、能缩并、能广播。但微分几何老师会告诉你：张量是坐标无关的几何对象，换基规则必须满足。这两种说法是冲突的同一件事，还是压根在谈两个不同层次的概念？

要回答这个问题，先退一步：我们最熟悉的"矩阵"，在什么意义上不只是一张数字表？

\subsection{从双线性型重新看矩阵}

设 \(V,W\) 为有限维实向量空间。一个双线性映射

\[
B: V\times W \to \R
\]

对两个输入分别线性。具体说，固定 \(w\) 时 \(B(\cdot,w)\) 是 \(V\) 上的线性函数；固定 \(v\) 时 \(B(v,\cdot)\) 是 \(W\) 上的线性函数。数学表达为：

\[
B(\alpha v_1+\beta v_2, w) = \alpha B(v_1,w) + \beta B(v_2,w),
\]

对第二个变量同样成立。

选好基。\(V\) 取 \(e_1,\ldots,e_m\)，\(W\) 取 \(f_1,\ldots,f_n\)。把所有基向量对的取值记下来：

\[
b_{ij} = B(e_i, f_j).
\]

这就得到了一个 \(m\times n\) 的数组 \([b_{ij}]\)——你熟悉的那个"矩阵"。任意输入

\[
v = \sum_i x_i e_i,\qquad w = \sum_j y_j f_j,
\]

代入多线性性质：

\[
B(v,w) = \sum_{i,j} b_{ij}\, x_i y_j = x^{\trans}[b_{ij}]y.
\]

矩阵 \([b_{ij}]\) 是双线性映射 \(B\) 在所选基下的坐标表示。没有什么神秘：数组是快照，双线性映射是被拍的对象。

换一组基会怎样？令 \(x=Sx'\)，则

\[
B(v,w) = x^{\trans}[b_{ij}]y = (x')^{\trans} S^{\trans}[b_{ij}]S y',
\]

所以新坐标下的矩阵是 \(S^{\trans}[b_{ij}]S\)。数字变了，双线性映射 \(B\) 没变。矩阵不是"一组数字"，而是"一个双线性映射在某组基下的数字快照，换了基要按照规则一起更新"。

这个视角把很多熟悉的东西统一起来。二次型 \(x^{\trans}Ax\) 是同一个向量塞进两个输入的双线性型。内积 \(\ip{x}{y}\) 是对称正定的双线性型。协方差 \(\Cov(X,Y)\) 是随机向量空间上的双线性型。它们都是矩阵，但矩阵是它们选了坐标之后的样子。

\subsection{从两个输入推广到多个输入}

把两个输入扩展成 \(k\) 个输入，就是高阶张量。给定向量空间 \(V_1,\ldots,V_k\)，一个 \(k\) 阶协变张量是多线性映射

\[
\mathcal{T}: V_1\times\cdots\times V_k \longrightarrow \R.
\]

它对每一个输入单独线性——无论是第几个槽，只要其他槽固定，这个槽上就是线性函数。

在第 \(r\) 个空间选基 \(e^{(r)}_1,\ldots,e^{(r)}_{n_r}\)。把所有基向量的 \(k\) 元组合上的取值记下来：

\[
t_{i_1\cdots i_k}
= \mathcal{T}\!\left(e^{(1)}_{i_1},\ldots,e^{(k)}_{i_k}\right).
\]

这些分量自然排成一个形状为

\[
n_1\times n_2\times\cdots\times n_k
\]

的多维数组。输入向量展开为坐标后，多线性性质立即给出：

\[
\mathcal{T}(v^{(1)},\ldots,v^{(k)})
= \sum_{i_1,\ldots,i_k}
t_{i_1\cdots i_k}\, x^{(1)}_{i_1}\cdots x^{(k)}_{i_k}.
\]

于是两个表述达成和解。``张量是多线性映射''说的是不依赖坐标的对象自身。``张量是多维数组''说的是选定基之后保存这个对象的分量表。前者回答"它是什么"，后者回答"怎样在机器里存"。工程上说的 tensor 通常指后者；数学里二者是同一枚硬币的两面。

工程文献还有一种用法：把向量空间张量积

\[
V_1\otimes\cdots\otimes V_k
\]

中的元素直接称为张量。在有限维 Euclidean 空间中，有内积自然地把向量和它的对偶（线性泛函）等同起来，\(V\) 和 \(V^*\) 之间的区分被内积消解了。没有内积，或进入微分几何的弯曲空间，协变（吃向量）和逆变（吃对偶向量，也就是被向量吃）的区分是必须严格对待的。本单元停在有限维实内积空间——这个前提下，你可以安全地把所有索引都当同等对待。

\subsection{阶与维度不是同一个数}

张量的阶（order）是索引方向的数量，也就是多线性映射的槽数。每个方向可以取多少个不同的坐标值，叫该模的维度（dimension）。两个概念完全独立：

\begin{itemize}
\tightlist
\item
  标量没有自由索引，是零阶张量；
\item
  向量有一个索引，是一阶张量；
\item
  矩阵有两个索引，是二阶张量；
\item
  \(T\times H\times W\times C\) 视频数组有四个索引，是四阶数据张量。
\end{itemize}

不要把阶与矩阵的秩（rank）搞混。一个三阶张量''阶为 3''只告诉你有三个模可以各自独立变换，没有告诉你存储它需要多少信息。三阶张量的 CP 秩可能为 5，也可能为 200——秩描述的是它能否拆成少数几个秩一成分之和。物理文献中常把 order 也称为 rank，跨界阅读时一定要先确认作者在说哪个数。

这个区分看起来只是术语问题，但写代码时很容易踩到。`np.ndim` 返回的是阶（几个轴），`np.linalg.matrix_rank` 返回的是矩阵秩（线性无关的行/列数）。对三阶数组，后者直接报错——因为它不知道该沿哪个模来做线性无关判断。

\subsection{秩一张量由外积产生}

矩阵中，秩一矩阵就是外积：\(uv^{\trans}\)，每个元素是 \(u_i v_j\)。高阶的直接推广：

给定向量 \(a^{(1)}\in\R^{n_1},\ldots,a^{(k)}\in\R^{n_k}\)，它们的外积

\[
\mathcal{X} = a^{(1)}\otimes\cdots\otimes a^{(k)}
\]

定义为

\[
x_{i_1\cdots i_k} = a^{(1)}_{i_1}\, a^{(2)}_{i_2}\cdots a^{(k)}_{i_k}.
\]

这就是秩一张量。每个模的变化由各自的向量独立控制，索引之间没有交叉耦合。三阶情形 \(x_{ijk}=a_i b_j c_k\) 可以读成：方向一的变化模式全在 \(a\) 里，方向二的全在 \(b\) 里，方向三的全在 \(c\) 里，三个方向的组合是纯粹乘出来的——没有任何 \(i\) 和 \(j\) 之间的"如果 \(i\) 大则 \(j\) 也大"那种交互。

矩阵 SVD 中的每一项 \(\sigma_r u_r v_r^{\trans}\) 正是二阶外积 \(\sigma_r(u_r\otimes v_r)\)。所以 SVD 已经是一种"拆成秩一成分之和"的做法。把它搬到高阶时，一个自然的期待是：我们也能把一张量拆成若干秩一张量的加权和。CP 分解（CANDECOMP/PARAFAC）正是这么做的。

但这里有一处令新手意外的地方。在矩阵 SVD 中，秩一成分的向量彼此正交（\(u_i\perp u_j\)，\(v_i\perp v_j\)），奇异值递减。这两种奢侈在高阶张量分解中通常会丢失：CP 分解的向量一般不正交，成分的贡献也不保证单调递减。下一篇进入具体分解方法时会展开。

\subsection{模乘法：只沿一个索引做线性变换}

矩阵乘法把两个二阶张量沿共同的索引缩并。但有时你只想沿张量的某一条轴施加线性变换，其余轴原样不动。这就是模乘法。

设三阶张量 \(\mathcal{X}\in\R^{I\times J\times K}\)，矩阵 \(M\in\R^{P\times I}\) 只作用在第一个模：

\[
\mathcal{Y} = \mathcal{X}\times_1 M,
\qquad
y_{pjk} = \sum_{i=1}^{I} m_{pi}\, x_{ijk}.
\]

第二、第三个索引 \(j,k\) 只是被原样搬运。类似地，\(\times_2 N\) 沿第二个模作用 \(N\)，\(\times_3 P\) 沿第三个模作用 \(P\)。多个模的变换可以串联，且互不干扰：

\[
\mathcal{X}\times_1 M \times_2 N \times_3 P.
\]

矩阵乘法只有两个槽：一个输入索引，一个输出索引。模乘法比矩阵乘法多了一层语义：它明确告诉你，线性变换挂在哪个轴上。

这个语义在工程中直接对应操作意图。视频的色通道沿第四模乘一个 \(3\times3\) 矩阵做颜色空间变换——RGB 转 YUV，每个像素独立处理，时间和空间位置不受影响。沿时间模乘 Fourier 或小波基，则只重组帧序列，颜色和空间不受影响。保留模的身份，使变换的物理含义可追踪。一旦你把多个模拍平成一个索引，再乘一个大矩阵，就很难向别人——甚至很难向自己——解释这一步变换到底改了什么。

\subsection{缩并：让一对索引相加消失}

矩阵乘法 \(C_{ik}=\sum_j A_{ij}B_{jk}\) 是一个经典的缩并操作：两个二阶张量各出一个索引（\(j\)），对位相乘后求和，结果还有两个自由索引（\(i,k\)）。内积 \(\ip{x}{y}=\sum_i x_i y_i\) 则是把两个一阶张量的唯一索引缩并掉，剩下一个标量。

对三阶张量，选择不同的索引做缩并产生不同的对象。沿第三模求和：

\[
M_{ij} = \sum_k x_{ijk},
\]

三阶变二阶——把所有"层"叠加成了一页。如果在求和时还带一个权重向量：

\[
M_{ij} = \sum_k x_{ijk} c_k,
\]

则是张量与向量沿第三模的缩并——每一层先乘上自己对应权重再加总。

缩并是几乎所有高阶运算的核心。Einstein 求和约定（同一项里上下重复出现的索引默认对它求和）在物理文献中把缩并写成极紧凑的形式，但初学时建议显式保留求和号。原因很简单：缩并掉的索引从结果中消失，自由索引保留在结果中。显式求和迫使你每次都要判断哪个索引在等号左边还出现，哪个不再出现——这是防止把自由索引与缩并索引混淆的最可靠检查。

一个实用的对照：PyTorch 的 `einsum('ijk,k->ij', X, c)` 等价于 \(M_{ij}=\sum_k x_{ijk}c_k\)。`einsum` 的字符串中，出现在箭头左边但不在右边的索引被缩并，同时出现在两边的保留为自由索引。这个设计就是把缩并的数学语义直接写进了接口。

\subsection{换基时，分量必须成组变化}

回到双线性型。换基后坐标矩阵变成 \(A'=S^{\trans}AS\)。这不是逐元素各自乘个系数——整张表要按照固定的模式联合更新。高阶张量同理。

三阶协变张量在三个基变换 \(S,R,U\) 下，分量按

\[
t'_{abc}
= \sum_{i,j,k} s_{ia}\, r_{jb}\, u_{kc}\, t_{ijk}
\]

联动变化。每个索引都按自己所属空间的基变换走一次。公式中具体的转置方向取决于基矩阵的约定（是旧基表新基还是新基表旧基），但不变量是明确的：每个索引都必须被自己的基变换矩阵独立处理。你只旋转数组的一条轴就声称整个几何对象没变——这通常是不完整的。

这正好划出了"多维数组"与"张量"之间的边界。数据处理中，人们常宽松地把任何多维数组叫张量——因为 NumPy 和 PyTorch 的类型系统只需要知道形状就能执行加法和乘法。数学上，一个数组要承担"张量"这个身份，还必须声明每个索引代表哪个空间（或哪个对偶空间），以及基变换时分量怎样联动。一个 \(3\times 3\times 3\) 数组可以是三阶应力张量，也可以是一个 mini-batch 的三个样本各 3 维特征——形状相同，换基规则完全不同。

因此，工程中它是 tensor 没问题——你的框架只需要形状。但当你把代码中的 tensor 解释为物理或几何量时，就需要回答换基规则。二者不矛盾；只是同一个词在不同的约束层级上使用。

\subsection{工程中的 shape 是语义约定，不只是尺寸检查}

两个数组可以有完全相同的 shape，却表示完全不同的对象。\(B\times T\times C\) 可以表示 batch、时间和通道（视频处理），也可以表示样本、节点和特征（图神经网络）。转置后元素总数不变——`np.transpose(X, (1,0,2))` 把 \(B\times T\times C\) 变成 \(T\times B\times C\)，尺寸检查通过，但原来的"第 \(b\) 个 batch 样本的第 \(t\) 帧"和"第 \(t\) 帧的第 \(b\) 个样本"是两件完全不同的事。后续的缩并操作——比如沿 batch 维求和求平均——在转置前的语义是"聚合所有样本"，在转置后变成了"聚合所有帧"。

因此，代码中的 reshape、transpose 和 broadcasting 每次操作时都应该能回答三个问题：哪些轴只是重新编号，哪些轴被混合，哪些索引将在下一步求和消失。shape 对上了只证明尺寸允许运算——编译器不替你验证语义。一个程序可以 shape 全对但完全算错事。

这就是为什么 einsum 式的显式索引记号比连续猜测 reshape 更不容易出错。`einsum('btc,tk->bck', X, W)` 明确说出了 b（batch）不参与缩并，t（时间）被缩并，c 和 k 是输出轴。意图落在字符串里，尺寸错误在运行时直接报 axis mismatch，不需要你心算 transpose 之后哪个轴变成了哪个。

\subsection{对称性减少独立分量}

若三阶张量满足

\[
t_{ijk} = t_{jik},
\]

它在前两个输入交换时不变，称在这两个指标上对称。若交换任意两个指标都不变，则是完全对称张量。三阶导数

\[
\frac{\partial^3 f}{\partial x_i\partial x_j\partial x_k}
\]

在足够光滑时，求导顺序可交换，所以对所有指标完全对称。

反对称张量在交换指标时变号：\(t_{ijk}=-t_{jik}\)。这种情况要求对角分量（两指标相等时）必然为零。行列式、有向体积和外代数中的楔积都与反对称结构相连。

对称性的实际影响体现在自由度上。一个 \(n\times n\times n\) 的完全对称三阶张量，独立分量数不是 \(n^3\)，而是 \(\binom{n+2}{3}\)。当对称性存在时，CP 分解或 Tucker 分解中的因子矩阵也应携带对应的对称约束——用无约束的分解去拟合对称数据，会浪费自由度并引入本不存在的"不对称噪声分量"。

从数据张量出发，后面两条路分叉了：下一篇考虑一般数据张量的 CP 与 Tucker 分解——怎样把"秩一成分之和"和"逐模子空间压缩"的想法从矩阵搬到高阶。若数据本身还有对称、非负、稀疏或物理守恒结构，分解时不应把这些当成事后约束，而应从一开始就嵌入分解假设。选什么分解，不仅是精度问题，也是你在声明数据背后的生成机制长什么样。
